\section{Pembagi Persekutuan Terbesar (PBB) dan Algoritma Euclidean}

Dalam kriptografi, kita sering perlu mengetahui apakah dua bilangan memiliki pembagi yang sama atau apakah sebuah bilangan memiliki invers modulo. Alat utama untuk melakukan hal ini adalah konsep \textbf{Pembagi Persekutuan Terbesar} (\textit{Greatest Common Divisor} atau GCD).

\subsection{Definisi PBB}

PBB dari dua bilangan bulat $a$ dan $b$ (tidak keduanya nol) adalah bilangan bulat positif terbesar yang membagi habis $a$ dan juga membagi habis $b$. Notasi yang digunakan adalah $\text{PBB}(a, b)$ atau $\gcd(a, b)$.

\textbf{Contoh}:
Pembagi dari 12 adalah $\{1, 2, 3, 4, 6, 12\}$.
Pembagi dari 18 adalah $\{1, 2, 3, 6, 9, 18\}$.
Pembagi persekutuan dari 12 dan 18 adalah $\{1, 2, 3, 6\}$.
Maka, $\text{PBB}(12, 18) = 6$.

\subsection{Bilangan Relatif Prima (\textit{Coprime})}

Dua bilangan bulat $a$ dan $b$ dikatakan \textbf{relatif prima} atau \textbf{coprime} jika pembagi persekutuan terbesar dari keduanya adalah 1:
\[ \text{PBB}(a, b) = 1 \]
Konsep ini sangat krusial dalam kriptografi. Sebagai contoh, dalam Affine Cipher, kunci perkalian $a$ harus relatif prima dengan modulus $n=26$. Jika $\text{PBB}(a, 26) \neq 1$, maka beberapa huruf plainteks akan menghasilkan cipherteks yang sama, sehingga pesan tidak dapat didekripsi secara unik.

\subsection{Algoritma Euclidean}

Untuk bilangan yang sangat besar, mencari PBB dengan mendaftar semua pembagi tidaklah efisien. Algoritma Euclidean menyediakan metode yang jauh lebih cepat berdasarkan prinsip bahwa PBB dari dua bilangan tidak berubah jika bilangan yang lebih besar digantikan oleh sisa baginya terhadap bilangan yang lebih kecil.

\textbf{Prinsip Dasar}:
\[ \text{PBB}(a, b) = \text{PBB}(b, a \pmod{b}) \]

\textbf{Langkah-langkah Algoritma Euclidean}:
\begin{enumerate}
    \item Misalkan kita ingin mencari $\text{PBB}(a, b)$ dengan $a > b$.
    \item Bagi $a$ dengan $b$ untuk mendapatkan sisa $r$: $a = qb + r$.
    \item Jika $r = 0$, maka $\text{PBB}(a, b) = b$.
    \item Jika $r \neq 0$, maka ulangi proses dengan mengganti $a$ menjadi $b$, dan $b$ menjadi $r$.
\end{enumerate}

\textbf{Contoh Numerik}: Hitung $\text{PBB}(252, 105)$.
\begin{align*}
    252 &= 2 \cdot 105 + 42 \quad \rightarrow \text{PBB}(252, 105) = \text{PBB}(105, 42) \\
    105 &= 2 \cdot 42 + 21 \quad \rightarrow \text{PBB}(105, 42) = \text{PBB}(42, 21) \\
    42 &= 2 \cdot 21 + 0 \quad \rightarrow \text{Selesai!}
\end{align*}
Maka, $\text{PBB}(252, 105) = 21$.

\subsection{Efisiensi Komputasi}

Kekuatan Algoritma Euclidean terletak pada kecepatannya. Secara matematis, jumlah langkah yang dibutuhkan untuk menemukan PBB tidak akan pernah melebihi lima kali jumlah digit dari bilangan yang lebih kecil (Teorema Lamé). Dalam notasi kompleksitas, algoritma ini berjalan dalam waktu $O(\log(\min(a, b)))$, yang membuatnya sangat efisien bahkan untuk angka dengan ribuan digit yang digunakan dalam standar keamanan modern.
